% GENERATED FILE. Edit canonical lessons, then run npm run generate:books.
% canonical-source-hash: fnv1a64:f89df3722950b1dc
% canonical-lessons: AR-C329-thawra, AR-C329-thara, AR-C329-istaqalla, AR-C329-istiqlal, AR-C329-shab

\chapter{Revolts and independence}
\label{ch:ar-revolts-and-independence}
\hlchaptermodality{329}

\begin{chapteropening}
\textbf{By the end of this chapter:} \emph{I can say a revolution, to revolt, to become independent, independence and a people.}

Complete the last lesson of chapter 329: I can say a revolution, to revolt, to become independent, independence and a people.
\end{chapteropening}

\section[\texorpdfstring{\ar{ثورة} (thawra)}{ثورة (thawra)}]{\ar{ثورة} (thawra) — a revolution, a revolt}

\label{lesson:AR-C329-thawra}



\begin{quote}
Before the new one: say the Arabic for a banquet, then the Arabic for an invited guest.
\end{quote}

\subsection*{You'll want to know: \ar{ثورة}}
\textbf{\ar{ثورة}} — \emph{thawra} — \textquotedblleft{}a revolution, a revolt\textquotedblright{}.

\begin{quote}
\ar{غيرت} \ar{الثورة} \ar{حياة} \ar{الناس}.
\end{quote}

\begin{quote}
\emph{ghayyarat ath-thawra ḥayāt an-nās.}
\end{quote}

\begin{quote}
The revolution changed people’s lives.
\end{quote}

\begin{cousinweb}
\textbf{\ar{ثورة}} (\emph{thawra}) is built on the native Arabic root \emph{th-w-r} (Lane), \textquotedblleft{}to rise up, to be stirred\textquotedblright{}, said of dust; the taught \emph{thawr}, a bull, is the same root.
\end{cousinweb}

\begin{grammarlens}[title={What you've built}]
One more word for telling what happened.
\end{grammarlens}

\subsection*{Your turn}
\begin{itemize}
\raggedright
  \item \emph{Say it:} \emph{thawra}
  \item \emph{Say it:} \emph{thawra}, once more
  \item \emph{Say it:} \emph{ghayyarat ath-thawra ḥayāt an-nās.} — The revolution changed people’s lives.
  \item \emph{Recall:} say the Arabic for a banquet, then the Arabic for an invited guest, then say \emph{thawra} again
\end{itemize}

\begin{tcolorbox}[breakable,colback=teal!4,colframe=teal!35!black,title={Before you move on}]
What does \ar{ثورة} mean? (\textquotedblleft{}A revolution, a revolt\textquotedblright{}.) Say it once more.
\end{tcolorbox}
\section[\texorpdfstring{\ar{ثار} (thāra)}{ثار (thāra)}]{\ar{ثار} (thāra) — to revolt; to flare up}

\label{lesson:AR-C329-thara}



\begin{quote}
Before the new one: say the Arabic for an invited guest, then the Arabic for a revolution.
\end{quote}

\subsection*{You'll want to know: \ar{ثار}}
\textbf{\ar{ثار}} — \emph{thāra} — \textquotedblleft{}to revolt; to flare up\textquotedblright{}.

\begin{quote}
\ar{ثار} \ar{الفلاحون} \ar{على} \ar{الحاكم} \ar{الظالم}.
\end{quote}

\begin{quote}
\emph{thāra al-fallāḥūn ʿalā al-ḥākim aẓ-ẓālim.}
\end{quote}

\begin{quote}
The farmers rose up against the unjust ruler.
\end{quote}

Present: \ar{يثور}, \emph{yathūr}.

\ar{ثار} \ar{غضبه}: his anger flared up.

\begin{cousinweb}
\textbf{\ar{ثار}} (\emph{thāra}) is built on the native Arabic root \emph{th-w-r} (Lane), \textquotedblleft{}to rise up\textquotedblright{}, said of dust and of anger, and then of a people.
\end{cousinweb}

\begin{grammarlens}[title={What you've built}]
One more word for telling what happened.
\end{grammarlens}

\subsection*{Your turn}
\begin{itemize}
\raggedright
  \item \emph{Say it:} \emph{thāra}
  \item \emph{Say it:} \emph{thāra}, once more
  \item \emph{Say it:} \emph{thāra al-fallāḥūn ʿalā al-ḥākim aẓ-ẓālim.} — The farmers rose up against the unjust ruler.
  \item \emph{Recall:} say the Arabic for an invited guest, then the Arabic for a revolution, then say \emph{thāra} again
\end{itemize}

\begin{tcolorbox}[breakable,colback=teal!4,colframe=teal!35!black,title={Before you move on}]
What does \ar{ثار} mean? (\textquotedblleft{}To revolt; to flare up\textquotedblright{}.) Say it once more.
\end{tcolorbox}
\section[\texorpdfstring{\ar{استقل} (istaqalla)}{استقل (istaqalla)}]{\ar{استقل} (istaqalla) — to become independent; to board}

\label{lesson:AR-C329-istaqalla}



\begin{quote}
Before the new one: say the Arabic for a revolution, then the Arabic for to revolt.
\end{quote}

\subsection*{You'll want to know: \ar{استقل}}
\textbf{\ar{استقل}} — \emph{istaqalla} — \textquotedblleft{}to become independent; to board\textquotedblright{}.

\begin{quote}
\ar{استقل} \ar{البلد} \ar{بعد} \ar{حرب} \ar{طويلة}.
\end{quote}

\begin{quote}
\emph{istaqalla al-balad baʿda ḥarb ṭawīla.}
\end{quote}

\begin{quote}
The country became independent after a long war.
\end{quote}

Present: \ar{يستقل}, \emph{yastaqill}.

It also takes a vehicle: \ar{استقل} \ar{القطار}, he took the train.

\begin{cousinweb}
\textbf{\ar{استقل}} (\emph{istaqalla}) is built on the native Arabic root \emph{q-l-l} (Lane), \textquotedblleft{}to be little; to lift\textquotedblright{}, the root of \emph{qalīl}, little, and of the taught \emph{aqall}, less. Form X \emph{istaqalla} is \textquotedblleft{}to bear oneself up\textquotedblright{}, alone.
\end{cousinweb}

\begin{grammarlens}[title={What you've built}]
One more word for telling what happened.
\end{grammarlens}

\subsection*{Your turn}
\begin{itemize}
\raggedright
  \item \emph{Say it:} \emph{istaqalla}
  \item \emph{Say it:} \emph{istaqalla}, once more
  \item \emph{Say it:} \emph{istaqalla al-balad baʿda ḥarb ṭawīla.} — The country became independent after a long war.
  \item \emph{Recall:} say the Arabic for a revolution, then the Arabic for to revolt, then say \emph{istaqalla} again
\end{itemize}

\begin{tcolorbox}[breakable,colback=teal!4,colframe=teal!35!black,title={Before you move on}]
What does \ar{استقل} mean? (\textquotedblleft{}To become independent; to board\textquotedblright{}.) Say it once more.
\end{tcolorbox}
\section[\texorpdfstring{\ar{استقلال} (istiqlāl)}{استقلال (istiqlāl)}]{\ar{استقلال} (istiqlāl) — independence}

\label{lesson:AR-C329-istiqlal}



\begin{quote}
Before the new one: say the Arabic for to revolt, then the Arabic for to become independent.
\end{quote}

\subsection*{You'll want to know: \ar{استقلال}}
\textbf{\ar{استقلال}} — \emph{istiqlāl} — \textquotedblleft{}independence\textquotedblright{}.

\begin{quote}
\ar{يحتفل} \ar{البلد} \ar{بعيد} \ar{الاستقلال} \ar{كل} \ar{سنة}.
\end{quote}

\begin{quote}
\emph{yaḥtafilu al-balad bi-ʿīd al-istiqlāl kull sana.}
\end{quote}

\begin{quote}
The country celebrates Independence Day every year.
\end{quote}

\begin{cousinweb}
\textbf{\ar{استقلال}} (\emph{istiqlāl}) is built on the native Arabic root \emph{q-l-l} (Lane), \textquotedblleft{}to lift\textquotedblright{}. \emph{Istiqlāl} is the verbal noun of Form X \emph{istaqalla}, \textquotedblleft{}to stand on one's own\textquotedblright{}.
\end{cousinweb}

\begin{grammarlens}[title={What you've built}]
One more word for telling what happened.
\end{grammarlens}

\subsection*{Your turn}
\begin{itemize}
\raggedright
  \item \emph{Say it:} \emph{istiqlāl}
  \item \emph{Say it:} \emph{istiqlāl}, once more
  \item \emph{Say it:} \emph{yaḥtafilu al-balad bi-ʿīd al-istiqlāl kull sana.} — The country celebrates Independence Day every year.
  \item \emph{Recall:} say the Arabic for to revolt, then the Arabic for to become independent, then say \emph{istiqlāl} again
\end{itemize}

\begin{tcolorbox}[breakable,colback=teal!4,colframe=teal!35!black,title={Before you move on}]
What does \ar{استقلال} mean? (\textquotedblleft{}Independence\textquotedblright{}.) Say it once more.
\end{tcolorbox}
\section[\texorpdfstring{\ar{شعب} (shaʿb)}{شعب (shaʿb)}]{\ar{شعب} (shaʿb) — a country’s people; a nation}

\label{lesson:AR-C329-shab}



\begin{quote}
Before the new one: say the Arabic for to become independent, then the Arabic for independence.
\end{quote}

\subsection*{You'll want to know: \ar{شعب}}
\textbf{\ar{شعب}} — \emph{shaʿb} — \textquotedblleft{}a country’s people; a nation\textquotedblright{}.

\begin{quote}
\ar{شعب} \ar{هذا} \ar{البلد} \ar{كريم} \ar{جدا}.
\end{quote}

\begin{quote}
\emph{shaʿb hādhā al-balad karīm jiddan.}
\end{quote}

\begin{quote}
The people of this country are very generous.
\end{quote}

Plural: \ar{شعوب}, \emph{shuʿūb}.

The adjective \ar{شعبي} is popular, folk: \ar{أكلة} \ar{شعبية}.

\begin{cousinweb}
\textbf{\ar{شعب}} (\emph{shaʿb}) is built on the native Arabic root \emph{sh-ʿ-b} (Lane), which means both \textquotedblleft{}to branch apart\textquotedblright{} and \textquotedblleft{}to gather\textquotedblright{}; Lane gives \emph{shaʿb} as a great tribe from which smaller ones branch.
\end{cousinweb}

\begin{grammarlens}[title={What you've built}]
One more word for telling what happened.
\end{grammarlens}

\subsection*{Your turn}
\begin{itemize}
\raggedright
  \item \emph{Say it:} \emph{shaʿb}
  \item \emph{Say it:} \emph{shaʿb}, once more
  \item \emph{Say it:} \emph{shaʿb hādhā al-balad karīm jiddan.} — The people of this country are very generous.
  \item \emph{Recall:} say the Arabic for to become independent, then the Arabic for independence, then say \emph{shaʿb} again
\end{itemize}

\begin{tcolorbox}[breakable,colback=teal!4,colframe=teal!35!black,title={Before you move on}]
What does \ar{شعب} mean? (\textquotedblleft{}A country’s people; a nation\textquotedblright{}.) Say it once more.
\end{tcolorbox}
